% =============================================================================
% style-small.tex —— 小版版式（本模板默认）
% -----------------------------------------------------------------------------
% 职责：只定义某一版「页面几何 + 文档字号 + 边距 + 页脚」这四件事。
%       公共配色、Block、页标题、过渡页等全部位于 config/ 下的其他文件，
%       因此切换版式不会影响观感，只会改变物理尺寸与字号。
%
% 加载时机：必须在 \documentclass{beamer} 之前 \input。
%   本文件用 \PassOptionsToClass 传字号，并在类加载后由 \SlideGeometry 设定纸张。
%
% 小版特征：16 × 9 cm（面积约为大版的 44%）、11pt；
%   页脚只有页码，没有横线和额外竖直留白。
% =============================================================================

% 文档字号：11pt（作用于所有相对字号，例如 \large、block 标题）
\PassOptionsToClass{11pt}{beamer}

% 页面尺寸：在 main.tex 于 \documentclass 之后立刻执行
\newcommand{\SlideGeometry}{\geometry{paperwidth=16cm,paperheight=9cm}}

% 左右边距：标题、上横线、正文共享同一数值，改这里即可整体收放
\newcommand{\InkMargin}{0.9cm}

% \sourceNote 的字号：小版用 \scriptsize
\newcommand{\SourceNoteSize}{\scriptsize}

% 版式相关字号与页脚：在导言区末尾由 main.tex 调用，此时公共模板已就绪
\newcommand{\ApplySlideStyle}{%
  \setbeamerfont{title}{size=\LARGE}
  \setbeamerfont{frametitle}{size=\fontsize{17}{21}}
  \setbeamerfont{section title}{size=\LARGE}
  \setbeamerfont{block title}{size=\fontsize{12}{15}}
  \setbeamerfont{subsection in toc}{size=\scriptsize}
  \setbeamerfont{footnote}{size=\scriptsize}
  \setbeamerfont{page number in head/foot}{size=\scriptsize}
  % 小版页脚：只有页码，无横线、无底部留白
  \setbeamertemplate{footline}{%
    \begin{beamercolorbox}[wd=\paperwidth,leftskip=\InkMargin,rightskip=\InkMargin]{page number in head/foot}
      \InkPageNumber
    \end{beamercolorbox}}}
